% !TEX program = lualatex
% ---------------------------------------------------------------------------
% Шаблон презентации курса «Методы анализа сетевых структур».
% Скопируйте папку template/ в новую папку рядом и правьте template.tex.
% Сборка: make -C .. template   (или latexmk -lualatex в этой папке, см. ../README.md)
% Здесь показаны все типовые слайды стиля.
% ---------------------------------------------------------------------------
\documentclass[aspectratio=169,10pt]{beamer}
\usetheme{HSEnet}
\usepackage[russian]{babel}
\usepackage{amsmath,amssymb,mathtools}
\usepackage{booktabs}
\usetikzlibrary{arrows.meta,positioning}

\course{Методы анализа сетевых структур}
\lecturelabel{Лекция N}
\title{Заголовок лекции в одну или две строки}
\subtitle{Подзаголовок: о чём лекция}
\author{Александр Пономаренко}
\institute{НИУ ВШЭ -- Нижний Новгород, магистратура «Интеллектуальный анализ данных»}
\date{Весна 2027}

\begin{document}
\maketitle

\section{Раздел лекции}

\begin{frame}{Обычный слайд: заголовок и список}
\begin{itemize}
\item Заголовок слайда --- утверждение или вопрос, а не тема из одного слова
\item Не больше 5--6 пунктов на слайде; пункт --- одна мысль
\item Акцент --- \alert{красным} или \textbf{полужирным}, не подчёркиванием
  \begin{itemize}
  \item второй уровень --- только когда без него никак
  \end{itemize}
\end{itemize}
\begin{block}{Блок: определение или формула}
\[
  L = D - A,\qquad x^\top L x = \sum_{(i,j)\in E} (x_i - x_j)^2 .
\]
\end{block}
\takeaway{Плашка-вывод: одна фраза, которую студент должен унести со слайда.}
\end{frame}

\begin{frame}{Две колонки: рисунок и текст}
\begin{columns}[T]
\column{0.5\textwidth}
\centering
\begin{tikzpicture}[v/.style={circle,draw=hseBlue,thick,fill=hsePale,minimum size=6mm,font=\small},
  e/.style={-{Latex[length=2mm]},thick,draw=hseBlue!70}]
  \node[v] (a) at (0,1.5) {$a$};  \node[v] (b) at (0,0) {$b$};
  \node[v] (c) at (2.5,0.75) {$c$}; \node[v,fill=hseOrange!30,draw=hseOrange] (d) at (4.5,0.75) {$d$};
  \draw[e] (a) -- (c); \draw[e] (b) -- (c); \draw[e,draw=hseOrange] (c) -- (d);
\end{tikzpicture}
\source{Источник рисунка или данных --- мелко и серым}
\column{0.5\textwidth}
\begin{alertblock}{Предупреждение}
Типичная ошибка или подводный камень.
\end{alertblock}
\begin{exampleblock}{Пример}
Конкретный числовой пример.
\end{exampleblock}
\end{columns}
\end{frame}

\begin{frame}{Таблица и палитра}
\begin{columns}[T]
\column{0.5\textwidth}
\small
\begin{tabular}{@{}lrr@{}}
\toprule
Метод & время & качество\\
\midrule
точный & 4 ч & \textbf{1{,}00}\\
эвристика & 10 мин & 0{,}97\\
\bottomrule
\end{tabular}
\column{0.5\textwidth}
\footnotesize
Цвета для рисунков --- только из палитры брендбука:\\[4pt]
\foreach \c in {hseBlue,hseLightBlue,hseSky,hseGreen,hseOrange,hseRed,hsePurple,hseYellow,hseSteel,hseTeal}{%
  \tikz\fill[\c] (0,0) circle (1.6mm);\,}
\\[4pt]
\foreach \c in {hsePale,hsePaleGreen,hsePaleRed,hsePaleYellow,hseLightGray,hseGray}{%
  \tikz\fill[\c] (0,0) circle (1.6mm);\,}
\end{columns}
\end{frame}

\begin{frame}[fragile]{Код}
\begin{semiverbatim}\small
import networkx as nx
G = nx.karate_club_graph()
L = nx.laplacian_matrix(G)
\end{semiverbatim}
\end{frame}

\finalframe{Вопросы?}{Александр Пономаренко\\ \href{https://github.com/aponom84/Network-Analysis-Course}{\color{white}github.com/aponom84/Network-Analysis-Course}}
\end{document}
